%%%%%%%%%%%%%%%%%%%%%%%%%%%%%%%%%%%%%%%%%%%%%%%%%%%%%%%%%%%%%%%%%%%%%%%%%%%%%%%%%%
\begin{frame}[fragile]\frametitle{}
\begin{center}
{\Large Dr Sushma Watwe यांच्या शिकवणी}
\end{center}

{\tiny (Ref: Dr Sushma Watwe --- ध्यान, YouTube)}

\end{frame}

%%%%%%%%%%%%%%%%%%%%%%%%%%%%%%%%%%%%%%%%%%%%%%%%%%%%%%%%%%%
\begin{frame}[fragile]\frametitle{ध्यान म्हणजे काय?}

	\begin{itemize}
	\item ध्यान = नो माइंड स्टेट --- मनात कोणताही विचार नसणे
	\item एकाग्रता (concentration) म्हणजे ध्यान नव्हे --- मन तिथे शिल्लक असते
	\item चिंतन (contemplation) म्हणजेही ध्यान नव्हे --- मन-बुद्धी दोन्ही सक्रिय
	\item ध्यानात: मन नाही, बुद्धी नाही --- पण `आपण आहोत' हे भान असते
	\item समाधीत ते भानही लुप्त होते
	\end{itemize}

{\tiny (Ref: Dr Sushma Watwe --- ध्यान)}

\end{frame}

%%%%%%%%%%%%%%%%%%%%%%%%%%%%%%%%%%%%%%%%%%%%%%%%%%%%%%%%%%%
\begin{frame}[fragile]\frametitle{सुशुप्ती विरुद्ध ध्यान}

	\begin{itemize}
	\item सुशुप्ती (गाढ झोप): कोणताच विचार नाही, मनाचे व्यापार नाही
	\item पण सुशुप्तीत अज्ञानाचे आवरण असते --- भान नसते
	\item ध्यानात मात्र जाणीव असते --- `काहीही नाही' हे सजाग अवस्थेत जाणते
	\item सुशुप्तीतील आनंद जागे झाल्यावर जाणवतो; ध्यानातील आनंद त्याक्षणी असतो
	\item लक्ष्य: सुशुप्तीसारखी निर्विचार स्थिती --- पण जागृत अवस्थेत
	\end{itemize}

{\tiny (Ref: Dr Sushma Watwe --- ध्यान)}

\end{frame}

%%%%%%%%%%%%%%%%%%%%%%%%%%%%%%%%%%%%%%%%%%%%%%%%%%%%%%%%%%%
\begin{frame}[fragile]\frametitle{ध्यानाचे फळ}

	\begin{itemize}
	\item ध्यानात मनाचा लय झाला की देह-विश्वाचे भान लुप्त होते
	\item पहिला टप्पा: \textbf{दुःख निवृत्ती} --- सर्व समस्या त्या क्षणापुरत्या लुप्त
		\begin{itemize}
		\item जसे आजारातून बरे वाटणे --- आजार गेला पण पूर्ण आरोग्य अजून नाही
		\end{itemize}
	\item दुसरा टप्पा: \textbf{निरतिशय आनंद} --- पुष्कळ साधनेनंतर
	\item अंतिम ध्येय: \textbf{आत्मदर्शन} --- `मी कोण आहे?' या प्रश्नाचे उत्तर
	\item रोज सारखेच ध्यान लागेल असे नाही --- देह-मन-बुद्धी-परिस्थितीवर अवलंबून
	\end{itemize}

{\tiny (Ref: Dr Sushma Watwe --- ध्यान)}

\end{frame}

%%%%%%%%%%%%%%%%%%%%%%%%%%%%%%%%%%%%%%%%%%%%%%%%%%%%%%%%%%%
\begin{frame}[fragile]\frametitle{ध्यानाची वेळ आणि स्थान}

	\begin{itemize}
	\item रात्री १२ ते ४ ही वेळ ध्यानासाठी टाळावी --- पवित्र मानली जात नाही
	\item झोप न लागल्यास: त्याच जागी नामस्मरण करत पडावे
	\item पाठांतर उपयुक्त --- मनाचे श्लोक, करुणाष्टक, गीता, दासबोध
	\item स्थान: शांत, पवित्र, घरातले ठराविक ठिकाण
	\item आसन: कडक (गाद्या नको), पाठ ताठ; दीड तास कम्फर्टेबली बसता येईल असे
	\item सद्गुरूंच्या सानिध्यात किंवा पवित्र वातावरणात ध्यान सहज लागते
	\end{itemize}

{\tiny (Ref: Dr Sushma Watwe --- ध्यान)}

\end{frame}

%%%%%%%%%%%%%%%%%%%%%%%%%%%%%%%%%%%%%%%%%%%%%%%%%%%%%%%%%%%
\begin{frame}[fragile]\frametitle{नामस्मरण --- ध्यानाचे प्रमुख साधन}

	\begin{itemize}
	\item चिंतनापेक्षा नामाचा उपयोग जास्त --- चिंतनाने मन गुंतते, नामाने सुटते
	\item नाम: इतर विचार थांबवण्यासाठी कवच
	\item नाम मनातल्या मनात, श्वासावर घ्यावे --- ओठ हलवायचे नाहीत
	\item नामात भक्तिभाव असल्यास त्याची कृपा मिळते
	\item सद्गुरूंचे, संतांचे, संप्रदायातील मंडळींचे स्मरण करून प्रार्थना करावी
	\item `मला ध्यान लागायला हवे' --- अशी प्रार्थना ते ऐकतात
	\end{itemize}

{\tiny (Ref: Dr Sushma Watwe --- ध्यान)}

\end{frame}

%%%%%%%%%%%%%%%%%%%%%%%%%%%%%%%%%%%%%%%%%%%%%%%%%%%%%%%%%%%
\begin{frame}[fragile]\frametitle{शक्ती संक्रमण}

	\begin{itemize}
	\item सद्गुरूंकडे शक्ती संक्रमणाची शक्ती असते --- अनुग्रहाच्या वेळी दिली जाते
	\item शक्ती ग्रहण करण्यास साधक पात्र असावा लागतो
	\item दगडावर बी टाकल्यास उपयोग नाही --- भूमी तयार हवी
	\item रामकृष्णाने विवेकानंदांना, रूपानंदाने माधवनाथांना शक्ती संक्रमित केली
	\item शक्ती संक्रमणानंतर साधनेचा पाठपुरावा न केल्यास ती नष्ट होते
	\item बी पेरल्यानंतर पाणी-पोषण न दिल्यास बी सुकते --- तसेच शक्तीचे
	\end{itemize}

{\tiny (Ref: Dr Sushma Watwe --- ध्यान)}

\end{frame}

%%%%%%%%%%%%%%%%%%%%%%%%%%%%%%%%%%%%%%%%%%%%%%%%%%%%%%%%%%%
\begin{frame}[fragile]\frametitle{यम-नियम आणि सात्विक आहार}

	\begin{itemize}
	\item यम-नियम पाळल्याशिवाय ध्यान लागणे शक्य नाही
	\item एकीकडे शुद्धीची प्रक्रिया, दुसरीकडे अशुद्ध खाणे --- विरोधाभास
	\item डॉक्टर औषध देतो, रुग्ण कुपथ्य करतो --- रोग कसा बरा होईल?
	\item सात्विक आहार आवश्यक --- तळकट, मसाले, बाहेरचे विपरीत
	\item ध्यानाने शरीर-मन-बुद्धीची शुद्धी होते --- आहाराने तीच टिकते
	\item जसजशी शक्ती जागृत होते तसतसे यमनियमांकडे अधिक लक्ष द्यावे
	\end{itemize}

{\tiny (Ref: Dr Sushma Watwe --- ध्यान)}

\end{frame}

%%%%%%%%%%%%%%%%%%%%%%%%%%%%%%%%%%%%%%%%%%%%%%%%%%%%%%%%%%%
\begin{frame}[fragile]\frametitle{इंद्रिय प्रत्याहार --- मोठा संघर्ष}

	\begin{itemize}
	\item ध्यानात सगळ्या बहिर्मुख वृत्ती अंतर्मुख करायच्या
	\item इंद्रियांना बाहेरच्या विषयाची गोडी --- हा मोठा संघर्ष
	\item संत सांगतात: ६-१२ वर्षे लागली मनाला आत आणायला
	\item बाहेरची संगती, खाणे-पिणे, वाचन, संवाद --- सर्व अनुकूल हवे
	\item संबंध नसलेल्या गोष्टींची चर्चा, वाचन, मेळावे टाळावे
	\item \textbf{नामाचे कवच:} प्रसंग आला तरी मन नामात ठेवा --- परिणाम होत नाही
	\end{itemize}

{\tiny (Ref: Dr Sushma Watwe --- ध्यान)}

\end{frame}

%%%%%%%%%%%%%%%%%%%%%%%%%%%%%%%%%%%%%%%%%%%%%%%%%%%%%%%%%%%
\begin{frame}[fragile]\frametitle{ध्यानाचे मार्ग --- मनोलय हे ध्येय}

	\begin{itemize}
	\item ध्यानाचे अनेक प्रकार: श्वासावर, नामावर, पाठांतर, शरीर-मन फिरवणे
	\item कुठल्याही मार्गाने जा --- \textbf{मन उन्मन करणे} हेच साधायचे
	\item भक्तीने भावावस्था येते, पण फार काळ टिकत नाही
	\item दीर्घकाळ टिकण्यासाठी \textbf{राजयोगाच्या} अंगाने जायला हवे
	\item शून्य $\rightarrow$ निशून्य $\rightarrow$ महाशून्य --- अशी वाटचाल
	\item संतांचे मन: विकार नाही, फक्त जगत् कल्याणाचा विचार
	\end{itemize}

{\tiny (Ref: Dr Sushma Watwe --- ध्यान)}

\end{frame}

%%%%%%%%%%%%%%%%%%%%%%%%%%%%%%%%%%%%%%%%%%%%%%%%%%%%%%%%%%%
\begin{frame}[fragile]\frametitle{समता दृष्टी}

	\begin{itemize}
	\item ध्यानाचा परिणाम: `सम सर्वेषु भूतेषु' --- सर्वत्र समत्व (गीता अ. ६)
	\item एकच चैतन्य सर्वत्र भरले आहे --- मुंगी-माशीपासून ब्रह्मदेवापर्यंत
	\item प्रकृतीकडे लक्ष नको, चैतन्याकडे पहा
	\item माझा आत्मा = त्याचा आत्मा --- देहाच्या पातळीवर भिन्नता, चैतन्यात एकच
	\item मला दुःख होते तसे दुसऱ्याला देऊ नये
	\item मीच विश्वरूप भरून राहिलो आहे --- त्यामुळे कुणाला दुःख देणे संभवत नाही
	\end{itemize}

{\tiny (Ref: Dr Sushma Watwe --- ध्यान)}

\end{frame}

%%%%%%%%%%%%%%%%%%%%%%%%%%%%%%%%%%%%%%%%%%%%%%%%%%%%%%%%%%%
\begin{frame}[fragile]\frametitle{जगमित्र भाव}

	\begin{itemize}
	\item भक्ताची व्याख्या: निंदा-स्तुती, मान-अपमान नाही (गीता अ. १२)
	\item अज्ञानी जीव आहे म्हणून त्रास देतो --- रागवायचे नाही
	\item संतांनी शत्रूबद्दलही कधी शत्रुत्व पत्करले नाही
	\item तुटे वाद-संवाद --- वाद घालण्यात अर्थ नाही; समजायचे असेल तर सांगणे
	\item `भगवंत त्याच्या रूपाने माझी कसोटी घेतो' --- अशी दृष्टी
	\item किती तास ध्यान केले यापेक्षा: आचार-विचारात हे येते का? हे तपासणे महत्त्वाचे
	\end{itemize}

{\tiny (Ref: Dr Sushma Watwe --- ध्यान)}

\end{frame}

%%%%%%%%%%%%%%%%%%%%%%%%%%%%%%%%%%%%%%%%%%%%%%%%%%%%%%%%%%%
\begin{frame}[fragile]\frametitle{कर्मयोग + भक्तियोग + ध्यान}

	\begin{itemize}
	\item ध्यान हे एकट्याने पुरेसे नाही --- तिन्ही एकत्र हवेत
	\item \textbf{श्रेयो ज्ञान अभ्यासात् ज्ञानात् ध्यान विशिष्यते}
	\item ध्यान म्हणजे ठराविक वेळेसाठीच नव्हे --- ज्ञानानंतर अखंड ध्यानाची स्थिती
	\item कर्म शुद्ध, निष्काम नसल्यास ध्यान खऱ्या अर्थाने लागत नाही
	\item ध्यानाला भक्तीची जोड द्यावी
	\item ध्यान लागले म्हणजे `मी ध्यान करत आहे' हे नसते --- सहज लागलेले ध्यान
	\end{itemize}

{\tiny (Ref: Dr Sushma Watwe --- ध्यान)}

\end{frame}

%%%%%%%%%%%%%%%%%%%%%%%%%%%%%%%%%%%%%%%%%%%%%%%%%%%%%%%%%%%
\begin{frame}[fragile]\frametitle{अकर्तृत्व बोध --- कर्मफल त्याग}

	\begin{itemize}
	\item ध्यानाने पूर्ण अलिप्तता: देह वेगळा, आत्मा वेगळा --- अनुभवात यायला हवे
	\item \textbf{कर्म जेथ दूरावे} (ज्ञानेश्वरी): मी अकर्ता --- कर्माशी संबंधच नाही
	\item दुसऱ्याने केलेल्या कर्माचे फळ मला नको --- कारण मी केलेच नाही
	\item तसे: देह-मन-बुद्धी कर्म करतात, आत्मा अकर्ता
	\item कर्तृत्व-भोक्तृत्व माझ्याकडे नसल्याने कर्मफलाची अपेक्षाच नाही
	\item परिणाम: शांती --- उदंड कर्म करूनही अत्यंत शांत राहणे
	\end{itemize}

{\tiny (Ref: Dr Sushma Watwe --- ध्यान)}

\end{frame}

%%%%%%%%%%%%%%%%%%%%%%%%%%%%%%%%%%%%%%%%%%%%%%%%%%%%%%%%%%%
\begin{frame}[fragile]\frametitle{सगुण + निर्गुण उपासना}

	\begin{itemize}
	\item ध्यानात दृश्य दिसणे म्हणजे: `मला ज्योत दिसली' --- अजून द्वैत आहे
	\item ती ज्योत पाहणाऱ्याकडे जाऊन बसायचे --- तोच मी आहे
	\item देवाची मूर्ती सगुणात पाहा, पूजा-अर्चा करा
	\item पण सगुण ज्या निर्गुणातून आले --- त्या निर्गुणाशी ध्यानात अनुसंधान साधायचे
	\item \textbf{दोन्ही एकत्र:} सगुण उपासना (नामस्मरण, पूजा) + निर्गुण (सोहम् बोध)
	\item हेच भगवंताला अभिप्रेत --- दोन्ही करणाऱ्यास ध्यानाचा पूर्ण उपयोग
	\end{itemize}

{\tiny (Ref: Dr Sushma Watwe --- ध्यान)}

\end{frame}

%%%%%%%%%%%%%%%%%%%%%%%%%%%%%%%%%%%%%%%%%%%%%%%%%%%%%%%%%%%
\begin{frame}[fragile]\frametitle{अंतिम ध्येय --- अहम् आत्मा}

	\begin{itemize}
	\item ध्यानाचे खरे ध्येय: देह ठीक राहण्यासाठी नव्हे --- आत्मदर्शनासाठी
	\item धारणा $\rightarrow$ ध्यान $\rightarrow$ समाधी --- या अंगाने वाटचाल
	\item \textbf{अहम् आत्मा} या बोधावर येऊन बसणे
	\item ज्ञान झाल्यानंतर: अज्ञानही नाही, सापेक्ष ज्ञानही नाही --- निंदा-स्तुती नाही
	\item वाईट विचार आले: चांगल्या संकल्पाने घालवणे
	\item चांगला संकल्प: `मला भगवत् प्राप्ती अपेक्षित आहे, माझे अवगुण जाऊ देत'
	\end{itemize}

{\tiny (Ref: Dr Sushma Watwe --- ध्यान)}

\end{frame}
